% Part: computability
% Chapter: recursive-functions
% Section: introduction

\documentclass[../../../include/open-logic-section]{subfiles}

\begin{document}

\olfileid{cmp}{rec}{int}
\olsection{تعارف}

قابلِ حسابیت دا ریاضیاتی نظریہ بنانے لئی سب توں پہلاں
قابلِ حسابیت دا اک \emph{نمونہ} بنانا پیندا اے۔ ہن اسی
قابلِ حسابیت نوں اوہ کم سمجھدے آں جو کمپیوٹر کردے نیں،
تے کمپیوٹر علامتاں نال کم کردے نیں۔ پر جدوں قابلِ حسابیت
دے نظریے بننے شروع ہوئے، حساب دی نمائندہ مثال \emph{عددی}
حساب سی۔ ریاضی دان ہمیشہ اعداد دے نظریے والے فنکشناں،
یعنی ایہو جیہے قابلِ حساب فنکشناں $f\colon \Nat^n \to
\Nat$، وچ دلچسپی رکھدے آئے نیں۔ ایس لئی حیرانی نہیں کہ
قابلِ حسابیت دے نظریے دے آغاز وچ انہاں ہی فنکشناں دا
مطالعہ ہویا۔ قابلِ حساب عددی فنکشناں دیاں سب توں جانے
پچھانے مثالاں، جیویں قدرتی عدداں اُتے جمع، ضرب تے قوت
چڑھاؤ، اک دلچسپ خاصیت سانجھی رکھدیاں نیں: انہاں نوں
\emph{بازگشتی طور اُتے} تعریف کیتا جا سکدا اے۔ ایس لئی
بازگشتی تعریفاں دے آسرے \emph{قابلِ حساب فنکشن} دی عام
تعریف دین دی کوشش فطری اے۔ اعداد دے نظریے والے فنکشناں
نوں بازگشتی طور اُتے تعریف کرن دے بہت سارے ممکنہ طریقیاں
وچوں اک خاصا سادہ تعریف دا نمونہ اتھے مرکزی بن جاندا اے:
جس نوں \emph{بنیادی بازگشت (primitive recursion)} آکھدے نیں۔

قابلِ حساب فنکشناں دے نال نال، اسی قابلِ حساب سیٹاں تے
تعلقاں وچ وی دلچسپی رکھ سکدے آں۔ کوئی سیٹ قابلِ حساب اے
جے اسی ایہ جواب حساب کر سکّیے کہ دتا ہویا عدد سیٹ دا
!!a{element} اے یا نہیں؛ تے کوئی تعلق اوہدوں تے صرف اوہدوں
قابلِ حساب اے جدوں اسی ایہ حساب کر سکّیے کہ ٹپل
$\tuple{n_1, \dots, n_k}$ تعلق دا !!a{element} اے یا نہیں۔
کسے سیٹ یا تعلق دے \emph{خصوصیتی فنکشن (characteristic
  function)} نوں ویکھ کے قابلِ حساب سیٹاں تے تعلقاں دی
بحث نوں قابلِ حساب فنکشناں دی بحث وچ سمویا جا سکدا اے۔
ایس طرح اسی بنیادی بازگشتی تعلق وی تعریف کر سکدے آں؛
مثال لئی تعلق ``$n$، $m$ نوں بغیر باقی دے تقسیم کردا اے''
اک بنیادی بازگشتی تعلق اے۔

بنیادی بازگشتی فنکشن---یعنی اوہ جنہاں نوں صرف بنیادی
بازگشت نال تعریف کیتا جا سکدا اے---پر اعداد دے نظریے
دے سارے قابلِ حساب فنکشن نہیں نیں۔ بنیادی بازگشت دیاں
کئی توسیعاں اُتے غور ہویا اے، پر فنکشناں نوں حساب کرن
دا سب توں طاقت ور تے بڑے پیمانے اُتے تسلیم شدہ اضافی
طریقہ بغیر مقررہ حد دی تلاش (unbounded search) اے۔ اِس
نال \emph{جزوی بازگشتی فنکشناں (partial recursive
  functions)} دی تعریف ملدی اے، تے اِس نال جڑی اک
تعریف \emph{عمومی بازگشتی فنکشناں (general recursive
  functions)} دی وی ملدی اے۔ عمومی بازگشتی فنکشن قابلِ
حساب تے کُل (total)، یعنی ہر ان پٹ اُتے تعریف شدہ، ہوندے
نیں؛ تے ایہ تعریف ٹھیک انہاں
جزوی بازگشتی فنکشناں دی نشان دہی کردی اے جو کُل ہون۔
بازگشتی فنکشن حساب دے ہر ہور نمونے (Turing machines،
lambda calculus وغیرہ) دی کارگزاری دی نقل کر سکدے نیں؛
ایس لئی اوہ حساب دے کئی تسلیم شدہ نمونیاں وچوں اک نیں۔

\end{document}
